\section{The precise arithmetic meaning of the result}
\label{sec:foundations}

We use first-order intuitionistic primitive recursive arithmetic, abbreviated
\ipra. Its language contains a symbol for every primitive recursive function,
with its defining equations; its logic is intuitionistic first-order logic
with equality, and its induction schema is restricted to quantifier-free
formulas. Natural numbers are the only objects. Integers, positive-denominator
rationals, Gaussian rationals, and finite lists are coded by natural numbers.
No quantification over a completed space of real or complex numbers is added.

The corresponding presentation in the inspected repository is
\texttt{PRAAx} in \texttt{ConciseIPRA/PR.lean}; despite its identifier, the
ambient deduction system is intuitionistic. Its induction interface is
\texttt{PRA.qfInd}. The theorem \texttt{pra\_iff\_native} transports a
\emph{derivation already constructed in this theory} to the native scalar
calculus. It does not turn an arbitrary theorem of ambient Lean into a
theorem of \ipra. All repository references in this package are pinned to
commit \texttt{b5a3ad037723c44d69a3d5b28a950c554327a989}; see
\cite{ProveIt} and the evidence manifest.

\subsection{Fixed algorithms, rather than a universal real-number sort}

An expression such as $\log u$, $e^z$, or $\zeta(s)$ below abbreviates
one specified family of rational computations, together with an error
bound. A finite assertion with requested error $2^{-v}$ is proved using
the finite computation at a prescribed finer precision. The notation for
the limiting value helps the reader see the mathematics; the arithmetic
claim is about those finite enclosures. In particular, we never postulate
a primitive recursive evaluator for all primitive recursive program codes.
The small, fixed collection of algorithms used here is compiled separately.

Bounded quantifiers cause no increase of the induction schema. For example,
if $R(i,\mathbf a)$ is a decidable primitive recursive predicate, then
``$R(i,\mathbf a)$ for all $i< B$'' is represented by the equation
\[
 \prod_{i<B}\chi_R(i,\mathbf a)=1.
\]
Products, sums, maximum, finite search with a supplied upper bound, and
the characteristic function in this equation are primitive recursive.
Induction applies to this single equation with all other parameters free.
The proof may universally generalize those parameters afterwards.

\subsection{A concrete logarithm algorithm}

For a positive rational $u\ge1$, find, by bounded integer search, the unique
$e\ge0$ with $2^e\le u<2^{e+1}$ and put $r=u/2^e$. Define
\[
 A_j(r)=2\sum_{i=0}^{j-1}\frac{z^{2i+1}}{2i+1},
 \qquad z=\frac{r-1}{r+1},\qquad 1\le r\le2.
\]
The finite geometric identity, integrated on $[0,z]$, gives
\[
 0\le \log r-A_j(r)
 =2\int_0^z\frac{t^{2j}}{1-t^2}\,dt
 \le\frac{2z^{2j+1}}{(2j+1)(1-z^2)}\le9^{-j}.
\]
This also proves convergence with an explicit bound, without invoking a
monotone-convergence principle. Set
\begin{equation}\label{eq:log-algorithm}
 \lambda(u,v)=eA_{v+e+2}(2)+A_{v+e+2}(r).
\end{equation}
Then
\begin{equation}\label{eq:log-error}
 0\le\log u-\lambda(u,v)\le(e+1)9^{-(v+e+2)}\le2^{-v}.
\end{equation}
The last inequality follows from $e+1\le2^{e+1}$ and $9\ge 2^3$.
For $u\ge2$, the initial term of $A_j(2)$ also gives
$\lambda(u,v)\ge2/3$. Rational inputs below $1$ are handled by
$\log u=-\log(1/u)$. Every loop has a bound supplied by the numerator,
denominator, and requested precision.

The ordinary logarithm identities used later can be proved directly for
these computations. One convenient common reference is
$\log u=\int_1^u dt/t$: the geometric identity above evaluates this
integral, substitution in finite Riemann sums proves
$\log(ab)=\log a+\log b$, and the derivative and monotonicity follow from
the explicit rational identity for $1/(u+h)-1/u$. The needed error bounds
on an interval with $u\ge a>0$ use only $a^{-1}$ and $a^{-2}$.
This gives finite rational proofs of these identities at every prescribed
precision; it does not appeal to an unspecified construction of real analysis.

\subsection{An explicit rational approximation to the offset integral}

We use the convention
\[
 \Li_2(x)=\int_2^x\frac{dt}{\log t}.
\]
It has no singular endpoint. Let $b(n)$ be the binary length of the
integer $n\ge2$, so $2^{b(n)-1}\le n<2^{b(n)}$. For $v\ge0$ set
\[
 h=2^{-(v+3)},\quad J=(n-2)2^{v+3},\quad
 P=v+b(n)+6,
\]
and define the nonnegative rational number
\begin{equation}\label{eq:Li-algorithm}
 \mathcal A(n,v)=h\sum_{j=0}^{J-1}
       \frac1{\lambda(2+jh,P)}.
\end{equation}
An empty sum, in particular at $n=2$, is zero. For $n<2$ we define the
algorithm to return zero, thereby making it total on all natural inputs.

Here is a complete error estimate. The function $f(t)=1/\log t$ is
positive and decreasing on $[2,n]$. The left and right Riemann sums on
this uniform mesh differ by $h(f(2)-f(n))\le2h$. Thus the exact left
sum exceeds the integral by at most $2h$. Since both logarithm
denominators are at least $2/3$, replacing $\log t$ by its lower
approximation increases its reciprocal by at most $3\cdot2^{-P}$.
Summing this error and using $n<2^{b(n)}$ gives
\begin{equation}\label{eq:Li-error}
 0\le\mathcal A(n,v)-\Li_2(n)
 \le 2h+3n2^{-P}
 <\left(\frac14+\frac3{64}\right)2^{-v}<2^{-v}.
\end{equation}
The executable reference implementation in the archive computes this
same rational sum. It is intended to specify a transparent algorithm,
not to provide an efficient way to evaluate the integral for very large $n$.

\subsection{The main first-order theorem}

Write $\Pi(n)$ for the number of integers $p\le n$ satisfying
$p\ge2$ and having no proper divisor between $2$ and $p-1$.
This definition uses only bounded divisibility tests and finite summation.
Set
\[
 C=2^{240},\qquad D=2C=2^{241},\qquad
 \mathcal N(m)=2^{D(m+7)^{10}}.
\]
The principal arithmetic conclusion of the argument is
\begin{theorem}[Finite logarithmic-integral strengthening]\label{thm:object}
In \ipra, for all natural $m,v,n$,
\begin{equation}\label{eq:object-target}
 n\ge\mathcal N(m)\quad\Longrightarrow\quad
 \left|\Pi(n)-\mathcal A(n,v)\right|
 \le \frac{n}{2^m b(n)}+2^{-v}.
\end{equation}
All functions and comparisons in this statement are primitive recursive.
After clearing the positive denominators, its matrix is quantifier-free.
\end{theorem}

For example, write $\mathcal A=a/d$ with $d>0$, without requiring the
fraction to be reduced, and let $z=|d\Pi(n)-a|$. The consequent is exactly
the integer comparison
\begin{equation}\label{eq:integer-matrix}
 z\,2^{m+v}b(n)
 \le d\bigl(n2^v+2^m b(n)\bigr).
\end{equation}
Thus the displayed theorem does not quantify over real numbers and does
not use an undecidable comparison between real-number codes.

The constants are chosen for simple verifiable estimates. They are
enormously larger than constants useful for numerical prime counting.
No claim of numerical optimization or a new classical error exponent is made.

\subsection{Elementary complex computations and their resource bounds}

For a Gaussian rational $z$ with $|z|\le B$, evaluate the exponential by
its Taylor polynomial. If $J$ is even and
$J\ge\max\{4\lceil B\rceil^2,2v+4\}$, then
\[
 \left|e^z-\sum_{j<J}\frac{z^j}{j!}\right|\le2^{-v}.
\]
Indeed, $J!\ge(J/2)^{J/2}$ gives $B^J/J!\le2^{-J/2}$.
Successive omitted terms then decrease by a factor at most $1/2$,
so their sum is at most $2^{1-J/2}\le2^{-v}$.
We always use the unambiguous safe choice
\begin{equation}\label{eq:exp-cutoff}
 J=4\lceil B\rceil^2+2v+4.
\end{equation}
Sine and cosine are the imaginary and real parts at $iz$; powers
$u^{-s}$ are evaluated as $e^{-s\log u}$. When the logarithm is also
approximated, a bound on the exponential derivative prescribes its
input precision. For $|z|\le B$ one may use the rational bound
$|e^z|\le3^{\lceil B\rceil+1}$, and ask for an input error at most
$2^{-v}/(3^{\lceil B\rceil+2}+1)$. Bit lengths turn this into a
natural-number precision by a bounded calculation.

Division is used only with a displayed lower bound. If $|w|\ge a>0$
and $|w-\widetilde w|\le\eta\le a/2$, then
\[
 |\widetilde w|\ge a/2,
 \qquad |w^{-1}-\widetilde w^{-1}|\le2\eta/a^2.
\]
The same estimate is obtained using squared norms and rational
arithmetic; square roots need not be decided. The zeta lower bounds
proved below are exactly what makes this operation legitimate.

For a real-parameter integral on a rational interval of length $\ell$,
a derivative bound $M$ makes a uniform mesh of width $h$ incur an error
at most $M\ell h$. A finite sum of $J$ terms is evaluated with
per-term error at most $2^{-v}/(J+1)$, adjusted for rational coefficients.
Products use their explicit magnitude bounds. These prescriptions give
precision bounds by addition, multiplication, exponentiation, and binary
length; no search for a successful accuracy is required.

There is a useful explicit rule for composing these estimates. Consider
the bounded table of $J$ rational arithmetic gates used in one calculation.
Let an integer $B\ge1$ bound the magnitudes of all exact intermediate
values and let every inverted denominator have modulus at least $1/B$.
Addition, multiplication, absolute value, and guarded inversion amplify
an input error at most $\epsilon\le1/(2B)$ by at most $8B^3$, including
one further local rounding error. Set the initial tolerance to
\[
 \epsilon_0=\frac{2^{-p}}{2B(8B^3)^{J+1}}.
\]
Quantifier-free induction on the number of evaluated rows proves that
the row-$j$ error is at most $(8B^3)^j\epsilon_0$ and that every guard
remains valid. The final error is less than $2^{-p}$. The bounds $J,B$
come from the displayed cutoffs and magnitude estimates. This is an
error-propagation lemma for a bounded arithmetic table with the specified
analytic macros, not a universal evaluator for primitive recursive codes.

\subsection{Why the proof uses only quantifier-free induction}

Every analytic operation used in the number-theoretic proof is reduced
below to these operations and two explicit tail estimates. The finite
contour lemma uses a predetermined rectangular grid, and its proof is
exact cancellation of internal edges plus a rational error sum. Euler
products are cut off by a supplied integer bound. The Perron integral
has a numerical truncation bound and an explicit finite scalar-kernel
calculation. No enumeration of zeta zeros, compactness of a continuum,
unbounded minimization, or principle of choice occurs.

Here is the induction discipline, which is essential to the claim of
Theorem~\ref{thm:object}. Fix the input numbers and the final accuracy.
Compute all resource bounds first. For a loop with counter $j\le J$,
the induction invariant is a finite list of rational identities and
inequalities for the first $j$ terms, with their accumulated error
allowances. Replace that list by the conjunction of its decidable
characteristic equations. Induct on this quantifier-free formula.
Induction never asserts ``the identity holds at every precision'' or
``every point of an infinite contour has the desired property''.
A claimed uniform analytic bound is instantiated at the finitely many
mesh points and at the finitely many interval bounds required by the
Taylor estimate.

For rational statements with positive slack, all apparent passages to
a limit are eliminated by choosing one sufficiently fine approximation
from the displayed bounds. For an exact identity in the expository
notation, use the finite version with an arbitrary prescribed error.
At the last step the spare margin in the prime-counting estimate and
\eqref{eq:Li-error} absorb these errors. Universal generalization over
the original inputs is made only after these finite proofs. This is a
description of the specific proof translation used here, not the false
general principle that every computable inequality is provable in \ipra.

\paragraph{Mechanization status.}
Theorem~\ref{thm:object} is established here at the level requested: a
human proof with its finite arithmetic interpretation. The archive does
\emph{not} contain a completed Lean proof term for it. The Lean files are
explicitly marked interfaces and outlines, and the final section lists
the proof objects still to be constructed. In particular, a successful
test of a rational evaluation program would not establish the universal
number-theoretic theorem.
